\chapter{What a Quantum Field Is}

These notes follow Leonard Susskind's third lecture in \emph{Particle Physics 1: Basic Concepts}, curated by LazyingArt LLC. The lecture resumes unfinished business from the previous meeting: periodic space, harmonic oscillators, and the relation between waves, quanta, and fields. The strategy is characteristically direct. We begin with a field on a circle, quantize the allowed wave numbers, label states by occupation numbers, introduce creation and annihilation operators mode by mode, and only then ask what the field itself is. From there the lecture turns to locality, Bose enhancement, stimulated emission, and the simplest algebra of bosonic fields.

\section{Periodic Space, Allowed Wave Numbers, and Momentum}

We do not begin with a field as an abstract operator. We begin with periodic space. The point of the circle is not metaphysical; it is a calculational trick. If we put a system in a finite interval with reflecting walls, momentum is spoiled by the reflections. If instead we impose periodic boundary conditions, a wave can continue around the circle without bouncing, and momentum conservation survives while the volume remains finite.

A plane wave on a circle of circumference \(L\) has the form \(e^{ikx}\). Periodicity requires that after going once around the circle the wave return to itself. Thus
\begin{equation}
e^{ikL}=1.
\end{equation}
This is the whole quantization condition. Since \(e^{i2\pi n}=1\) for any integer \(n\), we conclude that
\begin{equation}
kL=2\pi n,
\qquad
k=\frac{2\pi n}{L},
\qquad
n\in\mathbb{Z}.
\end{equation}
So the allowed wave numbers are discrete. This is not yet the number of particles. It is only the label of the allowed modes.

The same condition can be stated in the language of wavelength. If \(\lambda\) is the distance needed for one full cycle, then
\begin{equation}
k\lambda=2\pi,
\qquad
k=\frac{2\pi}{\lambda}.
\end{equation}
Large wavelength means small wave number, and small wavelength means large wave number. The lecture then reconnects this to ordinary quantum mechanics: the momentum of a particle described by such a wave is
\begin{equation}
p=\hbar k.
\end{equation}
In units with \(\hbar=1\), one often suppresses the distinction and speaks rather loosely of \(k\) as the momentum. For the present lecture that is harmless, provided we remember what has been quantized: the allowed values of \(k\), not the positions in space.

Later, when \(L\) is taken large, the spacing \(2\pi/L\) between neighboring values of \(k\) becomes small and the discrete sum over modes turns into an integral. But the lecture insists that the discrete version is the right place to begin, because it lets us see the bookkeeping cleanly.

\section{Occupation Numbers and Fock-Space State Labels}

Once the allowed momenta have been put in order, the next question is immediate: how many quanta occupy each allowed mode? This is where the lecture pauses over a genuine notational trap. There are now two different kinds of integers in the story. One integer labels the allowed values of \(k\) through \(k=2\pi n/L\). A different integer counts how many particles occupy a given mode. These must not be confused.

The second integer is the occupation number, written \(n(k)\). It means the number of quanta with wave number \(k\), or equivalently with momentum \(p=\hbar k\). A many-particle state is therefore specified by the entire list of occupations,
\begin{equation}
\big|\{n(k)\}\big\rangle
=
\big|\ldots,n(-2),n(-1),n(0),n(1),n(2),\ldots\big\rangle.
\end{equation}
The dots matter. Positive and negative wave numbers are both present, corresponding to right-moving and left-moving momentum states. In this language the vacuum is simply the state in which every slot is zero:
\begin{equation}
|0\rangle
=
|\ldots,0,0,0,0,0,\ldots\rangle.
\end{equation}

\begin{figure}[t]
\centering
\includegraphics[width=0.72\textwidth]{lecture_03_figure_02.png}
\caption{Occupation-number notation \(n(k)\) on the board. The screenshot is documentary evidence for the bookkeeping; the clean state notation is reconstructed in the text.}
\end{figure}

The lecture is deliberately concrete here. One does not yet need a Hamiltonian or a field equation. One only needs the idea that a state can be labeled by saying how many quanta sit in each momentum mode. This is the bridge from single-wave language to the many-particle Hilbert space.

\section{Mode Operators \(a^\dagger(k)\) and \(a(k)\)}

Now that the states have been labeled, the next step is to introduce operators that change those labels. On the board Susskind writes \(a^{+}(k)\) and \(a^{-}(k)\). In the notes we will use the standard notation \(a^\dagger(k)\) and \(a(k)\), but the meaning is the same: for each allowed momentum mode \(k\), there is one operator that raises the occupation number and one that lowers it.

Their action on an occupation-number state is
\begin{align}
a^\dagger(k)\big|\ldots,n(k),\ldots\big\rangle
&=
\sqrt{n(k)+1}\,\big|\ldots,n(k)+1,\ldots\big\rangle,
\\
a(k)\big|\ldots,n(k),\ldots\big\rangle
&=
\sqrt{n(k)}\,\big|\ldots,n(k)-1,\ldots\big\rangle.
\end{align}
The lecture momentarily stumbles over notation and then corrects itself in exactly the right way: what matters is that the operator acts only on the \(k\)-th slot, and that the square-root factor is part of the rule.

\begin{figure}[t]
\centering
\includegraphics[width=0.58\textwidth]{lecture_03_figure_03.png}
\caption{Creation and annihilation operators for a fixed mode \(k\). The board writes \(a^{+}(k)\) and \(a^{-}(k)\); the text adopts \(a^\dagger(k)\) and \(a(k)\).}
\end{figure}

At first these operators can look like elaborate bookkeeping devices. The lecture therefore immediately puts them to work. If we start from a one-particle state \(|k_7\rangle\), then a scattering process that removes the particle from mode \(k_7\) and places it in mode \(k_9\) is represented by
\begin{equation}
a^\dagger(k_9)\,a(k_7)\,|k_7\rangle = |k_9\rangle.
\end{equation}
In this example the square roots are trivial, because there is never more than one particle in the system.

The language already becomes more flexible when one allows the number of particles to change. A schematic \(1\to 2\) process may be represented by
\begin{equation}
a^\dagger(k_9)\,a^\dagger(k_{16})\,a(k_7)\,|k_7\rangle
=
|k_9,k_{16}\rangle.
\end{equation}
This could not be described in ordinary one-particle quantum mechanics. Here it is perfectly natural. But the lecture is careful: the existence of these operators does not mean that nature must create or destroy particles. It means that if the dynamics permits particle-number-changing processes, this is the language in which they can be described.

\section{From Classical Fourier Coefficients to the Quantum Field}

Only now does the lecture ask the central question: what is the field itself? The answer is given by stepping briefly back to the classical picture. A classical field on a periodic interval can be expanded in Fourier modes,
\begin{align}
\psi(x) &= \sum_k \alpha(k)e^{ikx},
\\
\psi^\ast(x) &= \sum_k \alpha^\ast(k)e^{-ikx}.
\end{align}
The coefficients \(\alpha(k)\) are the amplitudes of the Fourier components. In the classical theory they are ordinary complex numbers. Their real and imaginary parts are measurable amplitudes.

The quantum theory is obtained by replacing these Fourier coefficients with operators. The lecture briefly hesitates over which coefficient should be paired with creation and which with annihilation; this is only a convention, and the lecture explicitly corrects it on the fly. We will adopt the standard convention
\begin{equation}
\alpha(k)\rightsquigarrow a(k),
\qquad
\alpha^\ast(k)\rightsquigarrow a^\dagger(k).
\end{equation}
With that choice, the field becomes an operator-valued Fourier series:
\begin{align}
\Psi(x) &= \sum_k a(k)e^{ikx},
\\
\Psi^\dagger(x) &= \sum_k a^\dagger(k)e^{-ikx}.
\end{align}
Normalization factors can be chosen in different ways; the lecture keeps the formulas simple and postpones normalization refinements. We will do the same.

This is the major conceptual pivot of the lecture. The \(a\)'s are no longer just operators that act on the slots of a Fock-space state. They are the quantum-mechanical versions of the classical Fourier coefficients. The field is therefore not an ordinary function. It is an operator whose Fourier components create and annihilate quanta.

The lecture also stops here to head off a misunderstanding. Creation and annihilation operators are mathematical operators. To say that they can raise or lower particle number is not yet to say that the actual laws of nature do so. Whether particles are created, absorbed, emitted, or conserved is a dynamical question. But if such processes occur, these are the right variables.

\subsection{Question \& Answer}

What does \(\Psi^\dagger(x)\) actually create?

This is the sharpest local question in the lecture, and it is answered by direct computation. First, the annihilation field kills the vacuum:
\begin{equation}
\Psi(x)|0\rangle = 0,
\end{equation}
because there is no particle present to annihilate. The creation field is more interesting. Acting on the vacuum, it gives
\begin{align}
\Psi^\dagger(x)|0\rangle
&=
\sum_k e^{-ikx} a^\dagger(k)|0\rangle
\\
&=
\sum_k e^{-ikx}|k\rangle,
\end{align}
where
\begin{equation}
|k\rangle = a^\dagger(k)|0\rangle
\end{equation}
is the one-particle momentum eigenstate.

\begin{figure}[t]
\centering
\includegraphics[width=0.70\textwidth]{lecture_03_figure_05.png}
\caption{\(\Psi^\dagger(x)\) acting on the vacuum and being rewritten as a superposition of momentum eigenstates. The brace on the board marks the passage from operator language to state-vector language.}
\end{figure}

At first sight this seems to say that \(\Psi^\dagger(x)\) creates all momenta at once. But that is precisely how one-particle quantum mechanics builds a position state: one adds momentum eigenstates with phases \(e^{-ikx}\). So the lecture resolves the puzzle in the natural way:
\begin{equation}
\Psi^\dagger(x)\ \text{creates a particle localized at position }x.
\end{equation}
Equivalently, \(\Psi(x)\) annihilates a particle at position \(x\).

There is a nice side observation here. The coefficients \(e^{-ikx}\) all have unit magnitude, so this state has equal weight in momentum space. The field operator is therefore creating a coherent superposition of all allowed momenta, exactly tuned so that the resulting state is localized in position space.

\section{Local Reactions, Bose Enhancement, and Stimulated Emission}

Once \(\Psi(x)\) and \(\Psi^\dagger(x)\) have been given a positional meaning, the lecture turns immediately to locality. One could write fantastically nonlocal reaction operators if one wished: a particle is absorbed at one point and two particles reappear arbitrarily far away. But that is exactly what the lecture rejects. If we want a local theory, then the absorption and the emission should happen at the same point.

The local \(1\to 2\) process therefore has the schematic form
\begin{equation}
\Psi^\dagger(x)\Psi^\dagger(x)\Psi(x).
\end{equation}
A particle is removed at the point \(x\), and two particles are created at the same point.

\begin{figure}[t]
\centering
\includegraphics[width=0.92\textwidth]{lecture_03_figure_06.png}
\caption{Field expansion on the left board and the local probability-amplitude discussion on the right board. The screenshot preserves the lecture's actual board layout at the moment where the momentum-space calculation is built from the local field operators.}
\end{figure}

\begin{figure}[t]
\centering
\begin{tikzpicture}[scale=1.0]
\draw[->] (-2.2,0) -- (0,0);
\fill (0,0) circle (2.2pt);
\draw[->] (0,0) -- (2.0,1.1);
\draw[->] (0,0) -- (2.0,-1.1);
\node at (-1.2,0.28) {$k_{\mathrm{initial}}$};
\node[right] at (0.05,0.05) {$x$};
\node at (1.45,1.35) {$\ell$};
\node at (1.45,-1.35) {$m$};
\end{tikzpicture}
\caption{Editorial redraw of the local process emphasized in the lecture: one incoming particle is absorbed at the point \(x\), and two outgoing particles are emitted from that same point.}
\end{figure}

Let the initial state be a one-particle momentum eigenstate \(|k_{\mathrm{initial}}\rangle\). Then the first step is to let the annihilation field act:
\begin{align}
\Psi(x)|k_{\mathrm{initial}}\rangle
&=
\sum_q e^{iqx} a(q)a^\dagger(k_{\mathrm{initial}})|0\rangle
\\
&=
e^{ik_{\mathrm{initial}}x}|0\rangle.
\end{align}
Only the mode matching the occupied slot survives. The lecture is very explicit about this: an annihilation operator must find something to annihilate.

Now act with the two creation fields:
\begin{align}
\Psi^\dagger(x)\Psi^\dagger(x)\Psi(x)|k_{\mathrm{initial}}\rangle
&=
e^{ik_{\mathrm{initial}}x}
\sum_{\ell,m}
e^{-i\ell x}e^{-imx}
a^\dagger(\ell)a^\dagger(m)|0\rangle
\\
&=
\sum_{\ell,m}
e^{i(k_{\mathrm{initial}}-\ell-m)x}
|\ell,m\rangle.
\end{align}
Thus the final state is a superposition of all two-particle momentum states, with a pure phase as the coefficient. If one looks only at ordered pairs \((\ell,m)\), all such coefficients have the same magnitude.

The lecture then notices that one channel has been treated too quickly: the special case \(\ell=m\). If the first creation operator has already placed one particle in mode \(m\), then the second creation into that same mode does not come with coefficient \(1\), but with
\begin{equation}
\sqrt{2}.
\end{equation}
This is the first clear appearance of Bose enhancement. Creating into an already occupied mode is favored by the square-root rule.

Here the lecture becomes more interesting, not less, because a student points out a subtlety. When \(\ell\neq m\), the unordered two-particle final state may receive contributions from two ordered terms in the raw double sum, namely
\begin{equation}
a^\dagger(\ell)a^\dagger(m)|0\rangle
\qquad\text{and}\qquad
a^\dagger(m)a^\dagger(\ell)|0\rangle,
\end{equation}
and these are equal because creation operators commute. So the clean conclusion is not that every same-momentum final state must dominate every different-momentum one. The clean conclusion is narrower and more robust: whenever we create into an already occupied mode, an extra Bose factor appears.

The lecture then gives the clearer example. Suppose an atom emits one photon at the point \(x\). Acting on the vacuum,
\begin{equation}
\Psi^\dagger(x)|0\rangle
=
\sum_k e^{-ikx}|k\rangle,
\end{equation}
so all momenta appear with equal magnitude. That is spontaneous emission in this simplified model.

Now suppose there is already a photon in mode \(\ell\). Then acting with \(a^\dagger(k)\) on the existing state produces the usual result for \(k\neq \ell\), but for \(k=\ell\) one finds
\begin{equation}
a^\dagger(\ell)\big|\ldots,n_\ell,\ldots\big\rangle
=
\sqrt{n_\ell+1}\,
\big|\ldots,n_\ell+1,\ldots\big\rangle.
\end{equation}
Therefore emission into the already occupied mode is enhanced. The probability for that channel is multiplied by
\begin{equation}
n_\ell+1.
\end{equation}
If \(n_\ell=0\), this reduces to the spontaneous case. If \(n_\ell=100\), the probability is enhanced by a factor of \(101\). This is stimulated emission.

The lecture is careful about language here as well. Bosons do not literally attract one another into the same state. The enhancement is statistical. It comes from the algebra of identical particles, not from a new force.

\section{Bosonic Algebra, Number Density, and Late Clarifications}

Only after these examples does the lecture summarize the algebra. For bosons, creation operators commute with each other, annihilation operators commute with each other, and the only nontrivial commutator is
\begin{align}
[a^\dagger(k),a^\dagger(\ell)] &= 0,
\\
[a(k),a(\ell)] &= 0,
\\
[a(k),a^\dagger(\ell)] &= \delta_{k\ell}.
\end{align}
Here \(\delta_{k\ell}\) is the Kronecker delta: it is \(1\) when \(k=\ell\) and \(0\) otherwise. This is exactly the oscillator algebra, replicated mode by mode.

The lecture then asks for the meaning of the integrated operator
\begin{equation}
\int_0^L dx\,\Psi^\dagger(x)\Psi(x).
\end{equation}
Substituting the field expansions gives
\begin{align}
\int_0^L dx\,\Psi^\dagger(x)\Psi(x)
&=
\sum_{k,\ell}
a^\dagger(k)a(\ell)
\int_0^L dx\,e^{-ikx}e^{i\ell x}
\\
&=
\sum_{k,\ell}
a^\dagger(k)a(\ell)
\int_0^L dx\,e^{i(\ell-k)x}.
\end{align}
The integral over the circle is
\begin{equation}
\int_0^L dx\,e^{i(\ell-k)x}=L\,\delta_{k\ell},
\end{equation}
so the double sum collapses:
\begin{equation}
\int_0^L dx\,\Psi^\dagger(x)\Psi(x)
=
L\sum_k a^\dagger(k)a(k).
\end{equation}
But \(a^\dagger(k)a(k)\) is just the occupation number in mode \(k\). Therefore, up to the explicit normalization convention used in the lecture,
\begin{equation}
\frac{1}{L}\int_0^L dx\,\Psi^\dagger(x)\Psi(x)
=
\sum_k a^\dagger(k)a(k)
=
N,
\end{equation}
the total particle number. Correspondingly, \(\Psi^\dagger(x)\Psi(x)\) is the particle-density operator.

\subsection{Question \& Answer}

Why is the vacuum not the zero vector?

This question matters because earlier shorthand can make the two sound similar. The vacuum state \(|0\rangle\) is the normalized state with no particles in it. It is a perfectly good state:
\begin{equation}
\langle 0|0\rangle = 1.
\end{equation}
By contrast, the zero vector is the null vector of the Hilbert space. It does not describe any physical state at all.

So when one writes
\begin{equation}
a(k)|0\rangle = 0,
\end{equation}
the \(0\) on the right-hand side is the zero vector. The statement is that annihilating the vacuum produces no state. It is not the statement that the vacuum itself is the zero vector. The lecture insists on this distinction, and rightly so.

The late questions then widen the picture. In more than one spatial dimension, the periodic box becomes a torus. In two dimensions, one may think of a rectangle with opposite edges identified; in three dimensions, a cube with opposite faces identified. What becomes discrete is still momentum, now in each spatial direction. The periodic identification is a way of representing space so that one can work temporarily at finite volume and later let sums become integrals.

The lecture also returns to the analogy with the harmonic oscillator. The real and imaginary parts of the field play the role of conjugate oscillator variables. Up to conventional factors,
\begin{equation}
\Re \Psi \propto \Psi+\Psi^\dagger,
\qquad
\Im \Psi \propto \frac{\Psi-\Psi^\dagger}{i}.
\end{equation}
Because these are conjugate quantities, they do not commute. One cannot measure them simultaneously with arbitrary precision. But when the occupation numbers are very large, the commutator becomes small compared with the amplitudes themselves, and the field behaves approximately classically. That is why a large number of photons can be described, to a very good approximation, by an ordinary classical electromagnetic wave.

Finally, the lecture previews fermions. The key formal change is that commutators are replaced by anticommutators,
\begin{equation}
\{A,B\}=AB+BA.
\end{equation}
Schematically, the bosonic relation
\begin{equation}
[a(k),a^\dagger(\ell)]=\delta_{k\ell}
\end{equation}
is replaced in the fermionic case by an anticommutator relation of the same general type. The effect is exactly the opposite of Bose enhancement: the Pauli exclusion principle forbids two fermions from occupying the same state. In that sense, stimulated emission is a bosonic story.

\section{Summary}

The lecture has now built the simplest bosonic quantum field theory in the order in which one would actually want to discover it. Periodic space discretizes the allowed momenta. Occupation numbers label the many-particle states. Creation and annihilation operators change those occupations mode by mode. The quantum field itself is then defined as an operator-valued Fourier sum whose coefficients are precisely those creation and annihilation operators.

Once \(\Psi(x)\) and \(\Psi^\dagger(x)\) are understood as operators that annihilate and create particles at position \(x\), locality becomes the natural principle for writing reactions. From there the square-root factors in the oscillator algebra show their physical content: Bose enhancement and stimulated emission. Only after seeing that behavior do we summarize the algebra, identify \(\Psi^\dagger(x)\Psi(x)\) as the density operator, and note the late clarifications about the vacuum, higher-dimensional periodic spaces, the classical limit, and the coming contrast with fermions.